\documentclass[11pt,a4paper]{article}

% ------------------------------------------------------------
% Paquetes
% ------------------------------------------------------------
\usepackage[utf8]{inputenc}
\usepackage[T1]{fontenc}
\usepackage[spanish]{babel}
\usepackage{geometry}
\geometry{margin=2.5cm}
\usepackage{graphicx}      % Para insertar imagenes / capturas
\usepackage{listings}      % Para bloques de codigo
\usepackage{xcolor}
\usepackage{float}
\usepackage{hyperref}
\hypersetup{colorlinks=true, linkcolor=blue, urlcolor=blue}

% ------------------------------------------------------------
% Estilo para los bloques de codigo
% ------------------------------------------------------------
\definecolor{codegray}{rgb}{0.95,0.95,0.95}
\definecolor{codekeyword}{rgb}{0.0,0.0,0.6}
\definecolor{codestring}{rgb}{0.0,0.5,0.0}
\definecolor{codecomment}{rgb}{0.5,0.5,0.5}

\lstset{
  backgroundcolor=\color{codegray},
  basicstyle=\ttfamily\footnotesize,
  keywordstyle=\color{codekeyword}\bfseries,
  stringstyle=\color{codestring},
  commentstyle=\color{codecomment}\itshape,
  breaklines=true,
  frame=single,
  numbers=left,
  numberstyle=\tiny,
  showstringspaces=false,
  tabsize=2,
  captionpos=b
}

% ------------------------------------------------------------
% Portada
% ------------------------------------------------------------
\title{Conexión de Supabase con la aplicación Ionic/Angular}
\author{Desarrollo de Interfaces (0488) -- DAM2}
\date{\today}

\begin{document}

\maketitle

\section{Introducción}

En esta parte se documenta cómo se ha conectado el \textit{backend} de
\textbf{Supabase} con la aplicación desarrollada en \textbf{Ionic + Angular
(Standalone)}. La conexión se realiza mediante la librería oficial
\texttt{@supabase/supabase-js}, que crea un cliente a partir de la URL del
proyecto y la clave pública (\texttt{anon / publishable key}).

% ============================================================
% IMAGEN 1: Panel de Supabase "Connect to your project"
% ============================================================
% Captura del panel de Supabase con el boton "Connect" abierto,
% mostrando la opcion "Framework > Ionic Angular" y el comando de instalacion.
% Archivo sugerido: img/01-panel-connect.png
\begin{figure}[H]
  \centering
  % \includegraphics[width=0.9\textwidth]{img/01-panel-connect.png}
  \fbox{\parbox[c][5cm][c]{0.9\textwidth}{\centering
  [CAPTURA 1] Panel "Connect to your project" de Supabase.\\
  Mostrar la pestaña \textit{Framework} con \textbf{Ionic Angular}
  seleccionado y el comando \texttt{npm install @supabase/supabase-js}.}}
  \caption{Asistente de conexión de Supabase para Ionic Angular.}
  \label{fig:panel-connect}
\end{figure}

\section{Paso 1: Instalación de la librería}

El primer paso es instalar el SDK oficial de Supabase en el proyecto. Desde la
terminal integrada, en la raíz del proyecto, se ejecuta:

\begin{lstlisting}[language=bash,caption={Instalación del cliente de Supabase}]
npm install @supabase/supabase-js
\end{lstlisting}

Esto añade la dependencia al archivo \texttt{package.json} del proyecto.

% ============================================================
% IMAGEN 2: Terminal con el npm install correcto
% ============================================================
% Captura de la terminal de VS Code tras ejecutar el comando,
% mostrando el mensaje "added X packages".
% Archivo sugerido: img/02-npm-install.png
\begin{figure}[H]
  \centering
  % \includegraphics[width=0.9\textwidth]{img/02-npm-install.png}
  \fbox{\parbox[c][4cm][c]{0.9\textwidth}{\centering
  [CAPTURA 2] Terminal tras ejecutar \texttt{npm install @supabase/supabase-js},
  mostrando el mensaje \textit{"added ... packages"}.}}
  \caption{Resultado de la instalación de la dependencia.}
  \label{fig:npm-install}
\end{figure}

\section{Paso 2: Configuración de las credenciales (environments)}

Las credenciales del proyecto (URL y clave pública) se guardan en los archivos
de entorno de Angular. Así se separan de la lógica y se pueden cambiar entre
desarrollo y producción. La clave utilizada es la \textbf{publishable / anon
key}, que es pública y está pensada para el cliente.

\subsection*{src/environments/environment.ts}

\begin{lstlisting}[language=Java,caption={environment.ts (desarrollo)}]
export const environment = {
  production: false,
  supabaseUrl: 'https://rueljvohvzfdmtfjwazl.supabase.co',
  supabaseKey: 'sb_publishable_AhKlJL6ZZwxTRmuCQVOMSg_j2oCc86e',
};
\end{lstlisting}

\subsection*{src/environments/environment.prod.ts}

\begin{lstlisting}[language=Java,caption={environment.prod.ts (produccion)}]
export const environment = {
  production: true,
  supabaseUrl: 'https://rueljvohvzfdmtfjwazl.supabase.co',
  supabaseKey: 'sb_publishable_AhKlJL6ZZwxTRmuCQVOMSg_j2oCc86e',
};
\end{lstlisting}

% ============================================================
% IMAGEN 3: Donde copiar las credenciales en Supabase
% ============================================================
% Captura de Supabase en Project Settings > API (o del panel Connect)
% donde se ven la Project URL y la publishable key.
% Archivo sugerido: img/03-credenciales-supabase.png
\begin{figure}[H]
  \centering
  % \includegraphics[width=0.9\textwidth]{img/03-credenciales-supabase.png}
  \fbox{\parbox[c][5cm][c]{0.9\textwidth}{\centering
  [CAPTURA 3] Origen de las credenciales en Supabase:\\
  \textit{Project Settings > API} (o el panel \textit{Connect}),
  mostrando la \textbf{Project URL} y la \textbf{publishable key}.\\
  (No mostrar nunca la \texttt{service\_role key}).}}
  \caption{Credenciales del proyecto en el panel de Supabase.}
  \label{fig:credenciales}
\end{figure}

\section{Paso 3: Servicio de conexión (supabase.service.ts)}

Se crea un servicio de Angular (\texttt{providedIn: 'root'}, es decir, único
para toda la aplicación) que inicializa el cliente de Supabase con las
credenciales del \textit{environment}. Este servicio centraliza el acceso a la
base de datos.

\begin{lstlisting}[language=Java,caption={src/app/supabase.service.ts}]
import { Injectable } from '@angular/core';
import { createClient, SupabaseClient } from '@supabase/supabase-js';
import { environment } from '../environments/environment';

@Injectable({
  providedIn: 'root',
})
export class SupabaseService {
  private supabase: SupabaseClient;

  constructor() {
    this.supabase = createClient(
      environment.supabaseUrl,
      environment.supabaseKey
    );
  }

  getClientes() {
    return this.supabase.from('clientes').select('*');
  }
}
\end{lstlisting}

En el constructor, \texttt{createClient()} recibe la URL y la clave del
\textit{environment} y devuelve un objeto \texttt{SupabaseClient}. A partir de
ahí, cualquier consulta a la base de datos (por ejemplo \texttt{getClientes()})
se hace a través de ese cliente.

\section{Paso 4: Uso del servicio en una página}

Para comprobar la conexión, se inyecta el servicio en una página
(\texttt{Tab1}) y se realiza una consulta de prueba en \texttt{ngOnInit()}
sobre la tabla \texttt{clientes}.

\begin{lstlisting}[language=Java,caption={src/app/tab1/tab1.page.ts (fragmento)}]
export class Tab1Page implements OnInit {
  private supabaseService = inject(SupabaseService);

  clientes: any[] = [];
  error = '';

  async ngOnInit() {
    const { data, error } = await this.supabaseService.getClientes();
    if (error) {
      this.error = error.message;
      console.error('Error de Supabase:', error);
      return;
    }
    this.clientes = data ?? [];
    console.log('Clientes desde Supabase:', this.clientes);
  }
}
\end{lstlisting}

\section{Paso 5: Verificación de la conexión}

Para arrancar la aplicación en modo desarrollo se ejecuta:

\begin{lstlisting}[language=bash,caption={Arranque del servidor de desarrollo}]
npm start
\end{lstlisting}

La aplicación queda disponible en \url{http://localhost:4200/} (o en el puerto
que indique Ionic, p. ej. \texttt{8100}). Abriendo la consola del navegador
(tecla \texttt{F12}) se puede ver el resultado de la consulta, lo que confirma
que el cliente se comunica con Supabase. Se pueden dar dos resultados, ambos
válidos como prueba de conexión:
\begin{itemize}
  \item Si la tabla existe, se registra la lista de datos recibidos.
  \item Si la tabla aún no existe, el servidor de Supabase responde con un
        error \texttt{404} sobre dicha tabla. La petición igualmente llega al
        servidor correcto, lo que demuestra que la conexión está establecida;
        solo falta crear la tabla en la base de datos.
\end{itemize}

% ============================================================
% IMAGEN 4: App corriendo + consola del navegador
% ============================================================
% Captura de la app en localhost:4200 (pestana Tab1) junto con la
% consola del navegador (F12) mostrando el log "Todos desde Supabase: ...".
% Archivo sugerido: img/04-verificacion-consola.png
\begin{figure}[H]
  \centering
  % \includegraphics[width=0.9\textwidth]{img/04-verificacion-consola.png}
  \fbox{\parbox[c][5cm][c]{0.9\textwidth}{\centering
  [CAPTURA 4] Aplicación en el navegador (pestaña \textit{Tab 1 -- Todos
  (Supabase)}) con la consola (F12) abierta. La petición llega al servidor
  correcto (\texttt{...supabase.co/.../todos}), lo que confirma la conexión;
  el error \texttt{404} indica simplemente que la tabla \texttt{todos} aún
  no existe en la base de datos.}}
  \caption{Verificación de la conexión en el navegador.}
  \label{fig:verificacion}
\end{figure}

\section{Conclusión}

La conexión entre la aplicación Ionic/Angular y Supabase queda establecida a
través de:
\begin{enumerate}
  \item La librería \texttt{@supabase/supabase-js} instalada en el proyecto.
  \item Las credenciales (URL y clave pública) almacenadas en los
        \textit{environments}.
  \item Un servicio (\texttt{SupabaseService}) que crea el cliente y centraliza
        el acceso a la base de datos.
  \item Una consulta de prueba que confirma la comunicación con el servidor.
\end{enumerate}

\end{document}
